\begin{abstract}
Covariance regularisation is usually selected for estimation fidelity,
whereas zero-target cross-subject EEG is evaluated by discrimination on
unseen participants. We examine their mismatch along the complete scalar path
range, sampled using a predeclared finite grid in covariance/tangent-space linear discriminants. An operator
derivation shows that the standardized margin changes through
discriminant-path rotation while covariance predictive loss contains no
class-contrast term. Model selection and representation fitting used
neither target labels nor distribution statistics. On a prospectively
frozen EEGMMIDB block (\(N=20\)), source participant-macro-AUC selection
improved target macro AUC over source covariance-NLL selection by 0.0561
and Ledoit--Wolf by 0.0920 (both Holm-adjusted \(p\) values below 0.006).
The benefit is not uniform: 5/20 and 4/20 participants, respectively,
lost more than 0.02 AUC. In two independently frozen external cohorts
(BNCI2014\_002, \(N=14\); BNCI2015\_001, \(N=12\)), source objective
optima differed in 26/26 outer folds. Target effects were positive
(+0.0176 and +0.0433 versus covariance-NLL) but individually
non-significant after Holm correction (\(p=0.1397\) and 0.1856).
Their validity frontiers spanned 0.0384 and 0.0501 AUC with
cohort-dependent covariance-fidelity budgets. A separate prospective
experiment detected a regularisation-by-representation interaction
(-0.0778 AUC; Holm \(p=0.0133\)), opposite the historical filter-bank
hypothesis. We do not propose a universally superior estimator.
Instead, objective choice emerges as a consequential validity variable
in the studied cross-subject EEG pipelines.
\end{abstract}
